\medskip

\begin{minipage}{0.6\linewidth}
On considère un cube ABCDEFGH
et on appelle K le milieu du segment [BC].

On se place dans le repère 
$\left(\text{A}~;~\vect{\text{AB}},~\vect{\text{AD}},~\vect{\text{AE}}\right)$ 
et on considère  le tétraèdre EFGK.

On rappelle que le volume d'un tétraèdre est donné par: 

\[V = \dfrac13 \times \mathcal{B} \times h\]

où $\mathcal{B}$ désigne l'aire d'une base et $h$ la hauteur relative à cette base.
\end{minipage} \hfill
\begin{minipage}{0.35\linewidth}
\psset{unit=1cm}
\begin{pspicture}(5.5,5.8)
\psframe(0.2,0.2)(3.7,3.7)
\psline(3.7,0.2)(5,1.9)(5,5.4)(3.7,3.7)(4.35,1.05)
\psline(5,5.4)(1.5,5.4)(0.2,3.7)
\psline[linestyle=dashed](0.2,0.2)(1.5,1.9)(5,1.9)
\psline[linestyle=dashed](1.5,1.9)(1.5,5.4)
\pspolygon[linestyle=dotted,linewidth=1.25pt](0.2,3.7)(5,5.4)(4.35,1.05)
\uput[dl](0.2,0.2){\small A}\uput[dr](3.7,0.2){\small B}\uput[r](5,1.9){\small C}
\uput[dr](1.5,1.9){\small D}\uput[l](0.2,3.7){\small E}\uput[r](3.7,3.7){\small F}
\uput[ur](5,5.4){\small G}\uput[ul](1.5,5.4){\small H}\uput[dr](4.35,1.05){\small K}
\end{pspicture}
\end{minipage}

\medskip

\begin{enumerate}

\item Préciser les coordonnées des points E, F{}, G et K.
\item Montrer que le vecteur $\vect{n}\begin{pmatrix}\phantom{-}2\\-2\\\phantom{-}1\end{pmatrix}$ est orthogonal au plan (EGK).
\item Démontrer que le plan (EGK) admet pour équation cartésienne : $2x - 2y + z - 1 = 0.$
\item Déterminer une représentation paramétrique de la droite $(d)$ orthogonale au plan (EGK)
passant par F{}.
\item Montrer que le projeté orthogonal L de F sur le plan (EGK) a pour coordonnées $\left(\frac59~;~\frac49~;~\frac79\right)$.
\item Justifier que la longueur LF est égale à $\dfrac23$.
\item Calculer l'aire du triangle EFG. En déduire que le volume du tétraèdre EFGK est égal à $\dfrac16$.
\item Déduire des questions précédentes l'aire du triangle EGK.
\item On considère les points P milieu du segment [EG], M milieu du segment [EK] et N milieu du segment[GK]. Déterminer le volume du tétraèdre FPMN.
\end{enumerate}

\bigskip
